\documentclass[UTF8,12pt]{ctexart}
\usepackage[a4paper,margin=24mm]{geometry}
\usepackage{amsmath,amssymb,bm,graphicx,adjustbox}
\newcommand{\fitmath}[1]{\begin{adjustbox}{max width=\linewidth}$\displaystyle #1$\end{adjustbox}}
\setlength{\parindent}{0pt}
\setlength{\parskip}{.7em}
\sloppy
\allowdisplaybreaks
\begin{document}
\section*{1988 年考研数学三 · 真题}
\subsection*{1988年全国硕士研究生招生考试试题}

\subsection*{（试卷IV）}

\subsection*{一、填空题（本题共4小题，每小题3分，满分12分）}

（1）设 \(\displaystyle f(x)=\displaystyle \int_0^x e^{-t^2/2}\,dt\)，\(\displaystyle -\infty<x<+\infty\)，则：\textcircled{1} \(\displaystyle f\prime(x)=\)\_\_\_\_\_\_；\textcircled{2} \(\displaystyle f(x)\) 的单调性是\_\_\_\_\_\_；\textcircled{3} \(\displaystyle f(x)\) 的奇偶性是\_\_\_\_\_\_；\textcircled{4} 其图形的拐点是\_\_\_\_\_\_；\textcircled{5} 凹凸区间是\_\_\_\_\_\_；\textcircled{6} 水平渐近线是\_\_\_\_\_\_。


（2）行列式


\[\fitmath{\displaystyle \begin{vmatrix}1&1&1&0\\1&1&0&1\\1&0&1&1\\0&1&1&1\end{vmatrix}=\underline{\qquad}.}\]


（3）设矩阵


\[\fitmath{\displaystyle A=\begin{pmatrix}0&0&0&1\\0&0&1&0\\0&1&0&0\\1&0&0&0\end{pmatrix},\quad A^{-1}=\underline{\qquad}.}\]


（4）设 \(\displaystyle P(A)=0.4,\allowbreak \ P(A\cup B)=0.7\)，那么：\textcircled{1} 若 \(\displaystyle A\) 与 \(\displaystyle B\) 互不相容，则 \(\displaystyle P(B)=\)\_\_\_\_\_\_；\textcircled{2} 若 \(\displaystyle A\) 与 \(\displaystyle B\) 相互独立，则 \(\displaystyle P(B)=\)\_\_\_\_\_\_。


\subsection*{二、判断题（本题共5小题，每小题2分，满分10分）}

（1）若 \(\displaystyle \lim_{x\to x_0}f(x)\) 与 \(\displaystyle \lim_{x\to x_0}f(x)g(x)\) 均存在，则 \(\displaystyle \lim_{x\to x_0}g(x)\) 存在。


（2）若 \(\displaystyle x_0\) 是函数 \(\displaystyle f(x)\) 的极值点，则必有 \(\displaystyle f\prime(x_0)=0\)。


（3）等式 \(\displaystyle \int_0^a f(x)\,dx=-\displaystyle \int_0^a f(a-x)\,dx\) 对任何实数 \(\displaystyle a\) 都成立。


（4）若 \(\displaystyle A\) 和 \(\displaystyle B\) 都是 \(\displaystyle n\) 阶非零方阵，且 \(\displaystyle AB=O\)，则 \(\displaystyle A\) 的秩必小于 \(\displaystyle n\)。


（5）若事件 \(\displaystyle A,\allowbreak B,\allowbreak C\) 满足等式 \(\displaystyle A\cup C=B\cup C\)，则 \(\displaystyle A=B\)。


\subsection*{三、（本大题共4小题，每小题4分，满分16分）}

（1）求极限 \(\displaystyle \lim_{x\to1}\dfrac{\displaystyle x^x-1}{\displaystyle x\ln x}\)。


（2）已知 \(\displaystyle u+e^u=xy\)，求 \(\displaystyle \dfrac{\displaystyle \partial^2u}{\displaystyle \partial x\partial y}\)。


（3）求定积分 \(\displaystyle \int_0^3\dfrac{\displaystyle 1}{\displaystyle \sqrt x(1+x)}\,dx\)。


（4）求二重积分 \(\displaystyle \int_0^{\pi/6}dy\displaystyle \int_y^{\pi/6}\dfrac{\displaystyle \cos x}{\displaystyle x}\,dx\)。


\subsection*{四、（本大题共2小题，每小题3分，满分6分）}

（1）讨论级数 \(\displaystyle \sum_{n=1}^{\infty}\dfrac{\displaystyle (n+1)!}{\displaystyle n^{n+1}}\) 的敛散性。


（2）已知级数 \(\displaystyle \sum_{n=1}^{\infty}a_n^2\) 与 \(\displaystyle \sum_{n=1}^{\infty}b_n^2\) 都收敛，试证明级数 \(\displaystyle \sum_{n=1}^{\infty}a_nb_n\) 绝对收敛。


\subsection*{五、（本题满分8分）}

已知某商品的需求量 \(\displaystyle D\) 和供给量 \(\displaystyle S\) 都是价格 \(\displaystyle p\) 的函数：\(\displaystyle D=D(p)=\dfrac{\displaystyle a}{\displaystyle p^2},\allowbreak \ S=S(p)=bp\)，其中 \(\displaystyle a>0\) 和 \(\displaystyle b>0\) 为常数；价格 \(\displaystyle p\) 是时间 \(\displaystyle t\) 的函数且满足方程 \(\displaystyle \dfrac{\displaystyle dp}{\displaystyle dt}=k[D(p)-S(p)]\)（\(\displaystyle k\) 为正常数）。假设当 \(\displaystyle t=0\) 时价格为1，试求：


（1）需求量等于供给量时的均衡价格 \(\displaystyle p_e\)；（2）价格函数 \(\displaystyle p(t)\)；（3）极限 \(\displaystyle \lim_{t\to+\infty}p(t)\)。


\subsection*{六、（本题满分8分）}

在曲线 \(\displaystyle y=x^2\;(x\ge0)\) 上某点 \(\displaystyle A\) 处作一切线，使之与曲线以及 \(\displaystyle x\) 轴所围图形的面积为 \(\displaystyle \dfrac{\displaystyle 1}{\displaystyle 12}\)，试求：（1）切点 \(\displaystyle A\) 的坐标；（2）过切点 \(\displaystyle A\) 的切线方程；（3）由上述所围平面图形绕 \(\displaystyle x\) 轴旋转一周所成旋转体的体积。


\subsection*{七、（本题满分8分）}

已知线性方程组


\[\fitmath{\displaystyle \begin{cases}x_1+x_2+2x_3+3x_4=1,\\x_1+3x_2+6x_3+x_4=3,\\3x_1-x_2-k_1x_3+15x_4=3,\\x_1-5x_2-10x_3+12x_4=k_2.\end{cases}}\]


问 \(\displaystyle k_1\) 和 \(\displaystyle k_2\) 各取何值时，方程组无解？有唯一解？有无穷多解？在方程组有无穷多解的情况下，试求出一般解。


\subsection*{八、（本题满分7分）}

已知向量组 \(\displaystyle \boldsymbol\alpha_1,\allowbreak \boldsymbol\alpha_2,\allowbreak \ldots,\allowbreak \boldsymbol\alpha_s\;(s\ge2)\) 线性无关。设 \(\displaystyle \boldsymbol\beta_1=\boldsymbol\alpha_1+\boldsymbol\alpha_2,\allowbreak \ \boldsymbol\beta_2=\boldsymbol\alpha_2+\boldsymbol\alpha_3,\allowbreak \ldots,\allowbreak \boldsymbol\beta_{s-1}=\boldsymbol\alpha_{s-1}+\boldsymbol\alpha_s,\allowbreak \ \boldsymbol\beta_s=\boldsymbol\alpha_s+\boldsymbol\alpha_1\)。试讨论向量组 \(\displaystyle \boldsymbol\beta_1,\allowbreak \boldsymbol\beta_2,\allowbreak \ldots,\allowbreak \boldsymbol\beta_s\) 的线性相关性。


\subsection*{九、（本题满分6分）}

设 \(\displaystyle A\) 是3阶方阵，\(\displaystyle A^*\) 是 \(\displaystyle A\) 的伴随矩阵，\(\displaystyle A\) 的行列式 \(\displaystyle |A|=\dfrac{\displaystyle 1}{\displaystyle 2}\)，求行列式 \(\displaystyle |(3A)^{-1}-2A^*|\) 的值。


\subsection*{十、（本题满分7分）}

玻璃杯成箱出售，每箱20只。假设各箱含0，1，2只残次品的概率相应为0.8，0.1和0.1。一顾客欲购一箱玻璃杯，在购买时，售货员随意取一箱，而顾客随机地察看4只：若无残次品，则买下该箱玻璃杯，否则退回。试求：


（1）顾客买下该箱的概率 \(\displaystyle \alpha\)；（2）在顾客买下的一箱中，确实没有残次品的概率 \(\displaystyle \beta\)。


\subsection*{十一、（本题满分6分）}

某保险公司多年统计资料表明，在索赔户中被盗窃赔户占20\%，以 \(\displaystyle X\) 表示在随意抽查的100个索赔户中因被盗窃向保险公司索赔的户数。


（1）写出 \(\displaystyle X\) 的概率分布；（2）利用棣莫弗—拉普拉斯定理，求出被盗窃赔户不少于14户且不多于30户的概率的近似值。附表 \(\displaystyle \Phi(x)\) 是标准正态分布函数。


\[\fitmath{\displaystyle \begin{array}{c|ccccccc}x&0&0.5&1&1.5&2&2.5&3\\\hline\Phi(x)&0.5&0.692&0.841&0.933&0.977&0.994&0.999\end{array}}\]


\subsection*{十二、（本题满分6分）}

假设随机变量 \(\displaystyle X\) 在区间 \(\displaystyle (1,\allowbreak 2)\) 上服从均匀分布。试求随机变量 \(\displaystyle Y=e^{2X}\) 的概率密度 \(\displaystyle f(y)\)。


\subsection*{（试卷V）}

\subsection*{一、（本题满分12分）【同试卷IV 第一题】}

\subsection*{二、（本题满分10分）【同试卷IV 第二题】}

\subsection*{三、（本大题共4小题，每小题4分，满分16分）}

（1）求极限 \(\displaystyle \lim_{x\to1}(1-x^2)\tan\dfrac{\displaystyle \pi}{\displaystyle 2}x\)。


（2）已知 \(\displaystyle u=e^{x/y}\)，求 \(\displaystyle \dfrac{\displaystyle \partial^2u}{\displaystyle \partial x\partial y}\)。


（3）【同试卷IV 第三、（3）题】 （4）【同试卷IV 第三、（4）题】


\subsection*{四、（本题满分6分）}

确定常数 \(\displaystyle a\) 和 \(\displaystyle b\)，使函数


\[\fitmath{\displaystyle f(x)=\begin{cases}ax+b,&x>1,\\x^2,&x\le1\end{cases}}\]


处处可导。


\subsection*{五、（本题满分8分）}

将长为 \(\displaystyle a\) 的铁丝切成两段，一段围成正方形，另一段围成圆形。问这两段铁丝各长为多少时，正方形与圆形的面积之和最小？


\subsection*{六、（本题满分8分）【同试卷IV 第六题】}

\subsection*{七、（本题满分8分）【同试卷IV 第七题】}

\subsection*{八、（本题满分6分）}

已知 \(\displaystyle n\) 阶方阵 \(\displaystyle A\) 满足矩阵方程 \(\displaystyle A^2-3A-2E=O\)，其中 \(\displaystyle A\) 给定，\(\displaystyle E\) 是单位矩阵。证明：\(\displaystyle A\) 可逆，并求出其逆矩阵 \(\displaystyle A^{-1}\)。


\subsection*{九、（本题满分7分）【同试卷IV 第八题】}

\subsection*{十、（本题满分7分）【同试卷IV 第十题】}

\subsection*{十一、（本题满分7分）}

假设有十只同种电器元件，其中有两只废品。装配仪器时，从这批元件中任取一只，如是废品，则扔掉重新任取一只，直到取得正品之前，已取出的废品只数的概率分布、数学期望和方差。


\subsection*{十二、（本题满分5分）【同试卷IV 第十二题】}
\end{document}
